% Compilen desde esta carpeta o desde la raíz del repositorio.
\makeatletter
\def\input@path{{../../preamble/}{book2026-2027/preamble/}}
\makeatother
\input{preamble.tex}

\usepackage[spanish]{babel}
\usepackage{lmodern}
\usepackage{titlesec}
\geometry{a4paper,margin=1.6cm}
\linespread{1.10}

\definecolor{HeadingColor}{HTML}{243B53}
\definecolor{InstructionColor}{HTML}{704214}
\definecolor{ConceptColor}{HTML}{205C50}
\definecolor{LinkColor}{HTML}{0047AB}
\hypersetup{allcolors=LinkColor}
\newcommand{\guidelink}[2]{\href{#1}{\textbf{\underline{#2}}}}
\newcommand{\instruction}[1]{\textcolor{InstructionColor}{\textbf{#1}}}
\newcommand{\concept}[1]{\textcolor{ConceptColor}{\textbf{#1}}}
\newcommand{\checkbox}{%
    \begingroup
    \setlength{\fboxrule}{0.4pt}%
    \setlength{\fboxsep}{2pt}%
    \fbox{\rule{0pt}{0.8ex}\rule{0.8ex}{0pt}}%
    \endgroup}

\setlength{\parskip}{4pt}
\setlength{\tabcolsep}{5pt}
\renewcommand{\arraystretch}{1.15}
\setlist{leftmargin=0.65cm,topsep=4pt,partopsep=0pt,parsep=0pt,itemsep=3pt}
\titleformat{\section}{\large\bfseries\color{HeadingColor}}{}{0pt}{}
\titlespacing*{\section}{0pt}{10pt}{3pt}
\newcolumntype{Y}{>{\raggedright\arraybackslash}X}

\hypersetup{pdftitle={C1--C2. Coevalúen propuestas y fundamenten decisiones},pdfsubject={Emprendimiento e Innovación, Grado 10, Periodo 2}}

\begin{document}

\begin{center}
{\Large\bfseries\color{HeadingColor} C1--C2. Coevalúen propuestas y fundamenten decisiones\par}
\vspace{3pt}
Emprendimiento e Innovación \enspace|\enspace Grado 10 \enspace|\enspace Periodo 2\\
Guía 1 de 5
\end{center}

\section*{Objetivo y punto de partida}
\concept{C1.} Propone alternativas de solución variadas frente al reto de innovación definido.\\
\concept{C2.} Selecciona una solución según valor para el usuario, viabilidad y coherencia con el problema.

Recuperen el reto D1 (P1, C7) y los hallazgos H sobre el usuario (P1, C5). Participen en la evaluación de propuestas de otros equipos y usen lo aprendido para formular alternativas y justificar decisiones del proyecto propio.

\section*{Documento acumulativo e integridad de la evidencia}
Continúen el \textbf{documento acumulativo en LaTeX} en \guidelink{https://www.overleaf.com/}{Overleaf} o su editor habitual; entreguen fuente y PDF actualizado. Distingan \textbf{P1} y \textbf{P2}. Enlacen evidencias, decisiones, versiones y resultados por códigos o referencias internas, con fuentes y fechas reales.

Distingan \concept{observaciones y comentarios}, \concept{interpretaciones}, \concept{supuestos} e \concept{ideas}; identifiquen simulaciones y límites. \instruction{Usen IA solo para código LaTeX, formato y organización de contenido propio.} No inventen evidencias.

Actualicen la \textbf{bitácora} por sesión: fecha, aportes, evidencias, compromisos, responsables, plazos, estados y acuerdos. Las guías admiten proyectos físicos, digitales, comunicativos, de servicios o procedimentales. Adapten medios y extensión a la evidencia necesaria; los formatos sugeridos son opciones.

\section*{C1 de P2. Observen, pregunten y propongan alternativas}
\begin{enumerate}[series=pasos]
\item \textbf{Antes: preparen la coevaluación.} Lean el \guidelink{file:checklist_evaluativo_feria.pdf}{checklist de coevaluación}. Acuerden quién observa, pregunta y registra, con participación de todo el equipo. Identifiquen grupo evaluador, curso y proyectos evaluados. Usen los espacios según las propuestas observadas, sin inventar datos para llenar casillas.
\item \textbf{Durante: escuchen y documenten.} Observen la explicación y los recursos de cada equipo. Formulen preguntas para aclarar usuario, necesidad, funcionamiento, respaldo o viabilidad. Registren preguntas, respuestas y hechos localizables por proyecto; separen lo visto u oído de su interpretación. Eviten juicios como «es excelente» sin indicar qué los respalda.
\item \textbf{Evalúen con el instrumento intacto.} Atiendan sus siete categorías: comprensión de la problemática; claridad y organización de la información; dominio y conocimiento del tema; innovación y pertinencia de la solución; viabilidad e impacto; comunicación y uso de recursos; gestión del tiempo y trabajo en equipo. En cada ítem usen \textbf{E = Se evidencia, P = Parcialmente, A = Aún no se evidencia, N = No se observa}. Vinculen las marcas a sus notas; si algo no se observó, no lo conviertan en una afirmación de inexistencia. Estas categorías son del checklist: \textbf{no equivalen a C1--C10 de P2}.
\item \textbf{Después: elaboren alternativas.} Completen por proyecto los campos \textbf{Problemática, Fortaleza, Mejora concreta, Solución alternativa y Cómo o por qué}. Expliquen cómo funcionaría la alternativa o por qué ayudaría. Desarrollen alternativas con mecanismos distintos, a partir de lo observado, las preguntas y las ideas propias; identifiquen a qué problemática responde cada una y su origen. No basta cambiar el nombre o la apariencia de una idea.
\end{enumerate}

\section*{Entregables de C1}
\begin{itemize}
\item \textbf{Checklist legible}, con marcas y campos completos para los proyectos observados y notas que los respalden; una marca sola no demuestra C1.
\item \textbf{Solución alternativa y Cómo o por qué}, con elaboración propia que muestre variedad de mecanismos, origen y relación con las problemáticas: evidencia específica de C1.
\end{itemize}

\newpage
\section*{C2 de P2. Comparen y justifiquen una selección}
\begin{enumerate}[resume=pasos]
\item \textbf{Analicen la retroalimentación.} Revisen las marcas y los campos Problemática, Fortaleza y Mejora concreta junto con sus respaldos. Compartan respetuosamente las devoluciones con los equipos evaluados; registren las aclaraciones recibidas y qué cambió en su interpretación. Si reciben comentarios sobre su proyecto, identifiquen su procedencia. No atribuyan comentarios que no recibieron.
\item \textbf{Conecten con el proyecto propio.} Expliquen qué aprendizajes de la coevaluación son pertinentes para su usuario y D1, y cuáles no pueden trasladar sin comprobarlos. Retomen sus ideas de P1 y desarrollen alternativas propias (por ejemplo, ALT1, ALT2). Citen las propuestas ajenas que inspiraron una idea y expliquen la adaptación, sin presentarlas como creación del equipo.
\item \textbf{Comparen opciones con el mismo marco.} Organicen una comparación que identifique cada alternativa y su mecanismo, el valor esperado para el usuario y su respaldo, la viabilidad (recursos, tiempo, acceso y condiciones) y la coherencia con D1. Pueden usar una matriz o un análisis por opciones. Marquen «sin datos» cuando falte evidencia y señalen cómo obtenerla; las marcas del checklist no sustituyen esta comparación.
\item \textbf{Justifiquen la selección.} Expliquen qué opción eligen para el proyecto propio, frente a las opciones descartadas, y qué evidencias o condiciones pesaron en la decisión. Distingan la evidencia del usuario de lo aprendido al observar otros proyectos. Si dependen de un supuesto decisivo, declaren la elección provisional y qué falta comprobar. Registren el acuerdo y los compromisos en la bitácora.
\item \textbf{Sinteticen la sección de P1 del documento acumulativo.} Déjenla lo más concisa posible: conserven solo lo más relevante para comprender necesidad, usuario, reto definido, hallazgos que lo sustentan y decisiones que orientan el proyecto actual. Eliminen redundancias, integren información repetida y sinteticen descripciones que ya no aportan a esa comprensión. Preserven la redacción oficial de los criterios y las referencias necesarias. Conserven evidencias originales y registros completos en anexos o archivos de respaldo identificados, accesibles y referenciados desde la síntesis para mantener la trazabilidad. Organicen P1 sin rehacer el trabajo.
\end{enumerate}

\section*{Entregables de C2}
\begin{itemize}
\item \textbf{Análisis del checklist}: marcas respaldadas y campos \textbf{Problemática, Fortaleza y Mejora concreta}, relacionados con las alternativas y referenciados a observaciones, preguntas y comentarios. Expliquen qué aprendizaje adoptan y qué decisión justifican con él.
\item \textbf{Comparación propia y selección argumentada}: alternativas, valor para el usuario, viabilidad y coherencia con D1, fuentes, condiciones y razones para elegir y descartar opciones. Distingan evidencia, supuestos y comprobaciones pendientes. \textbf{La comparación y la decisión propias evidencian C2}; el instrumento aporta insumos.
\item \textbf{Sección de P1 sintetizada}, con criterios oficiales intactos y referencias accesibles a las evidencias originales y registros completos. Esta organización del documento \textbf{no sustituye la evidencia de selección justificada de C2}.
\end{itemize}

\section*{Resultado y continuidad}
Conserven instrumento y notas en coevaluación; alternativas, elección y supuestos en el proyecto propio. C3--C4 retomarán la síntesis de P1, sus respaldos y la selección de C2 para representar y materializar la elección. La bitácora respaldará la revisión de la colaboración en C9.

\section*{Pregunta de reflexión}
¿Cómo influyó la coevaluación en su elección de solución, a la luz de las evidencias del usuario y las incertidumbres reconocidas al comparar alternativas?

Respondan en el documento acumulativo con apoyo en su experiencia y evidencias como equipo.

\end{document}
